\begin{frame}{Проблема 5: частичные и неоднозначные сбои}
	\begin{itemize}
		\item Для fault-tolerant RTOS недостаточно учитывать только fail-stop.
		\item После сбоя узлы могут иметь неактуальное представление о состоянии друг друга.
		\item Для группы вычислителей нужны:
		\begin{itemize}
			\item heartbeat;
			\item local view;
			\item suspicion level;
			\item majority voting;
			\item controlled rejoin после recovery.
		\end{itemize}
	\end{itemize}
	
	\vspace{0.3cm}
	
	\begin{block}{Роль опорной статьи}
		В статье imperfect knowledge рассматривается как важная часть recovery-модели.
		Это делает её особенно полезной для fault-tolerant RTOS.
	\end{block}
\end{frame}